\documentclass[12pt,a4paper,twoside]{report}

% Wybierz styl nagłówka rozdziału:
% 'classic' — Rozdział X (nad tytułem)
% 'inline'  — X. Tytuł rozdziału (w jednej linii)

\def\chapterstyletype{classic}  % lub zmień na 'inline'

%MARGINESY
\usepackage[left=3cm, right=2.5cm, top=2.5cm, bottom=2.5cm, headsep=1.2cm]{geometry}

%INTERLINIA
\linespread{1.3}

%Język polski
\usepackage{polski}
\usepackage[utf8]{inputenc}
\usepackage[T1]{fontenc}
\usepackage{longtable}

% Pakiet, który poprawia typografię: kerning, justowanie i drobne szczegóły tekstu
\usepackage{microtype}

%landscape obrocone
\usepackage{pdflscape}

%tekst matematyczny
\usepackage{amsmath}
\usepackage{amsfonts}
\usepackage{amssymb}
\usepackage{amsthm}

%dodawanie grafiki
\usepackage{graphicx}

%do grupowania rysunków
\usepackage{subcaption}

%brak wdów - pojedyncze linie
\usepackage[all]{nowidow}

%Zapobieganie dzieleniu się wyrazów
%\tolerance=1
\emergencystretch=\maxdimen
\hyphenpenalty=10000
\hbadness=10000

%do kolorów
%\usepackage{color}
\usepackage{xcolor}

\usepackage{textcomp}%do symbolu euro
\usepackage[official]{eurosym}%drugi symbol euro

\usepackage[colorlinks=true,allcolors=black]{hyperref}

\usepackage{indentfirst} %dodaje wcięcie w pierwszym akapicie

%listy numerowane
\usepackage{enumitem}
\setlist[enumerate]{noitemsep,nolistsep}
\setlist[itemize]{noitemsep,nolistsep} 

%do modyfikacji wyglądu spisu treści, listy rysunków i tabel
\usepackage{tocloft}

%do tabel
\usepackage{tabularx}
\usepackage{multirow} %do połączenia wierszy
\usepackage{diagbox}%do ukośnej linii w tabelii
\usepackage{booktabs}
\usepackage{array}
\newcolumntype{Y}{>{\centering\arraybackslash}X}
\newcolumntype{j}{>{\raggedleft}p{200pt}}
\usepackage{caption}
\usepackage{lscape}





%podpisy rysunków i tabel
\renewcommand*{\figurename}{Rys.} %zmiana podpisu rysunku
\renewcommand*{\tablename}{Tab.} %zmiana podpisu tabeli
\captionsetup{font=small, labelsep=period, justification=centerlast} % Modyfikacja stylu podpisów: 
% - font=small ustawia rozmiar czcionki podpisu na 11 pt (można użyć footnotesize dla 10 pt),
% - labelsep=period wprowadza kropkę po numerze (zamiast dwukropka),
% - justification=centerlast centrowanie podpisu z wyrównaniem ostatniego wiersza do środka.


%%%%%%%%%%%%%%%%%%%%%%%%%%%%%%%%%%%%%%%%%%%%%%%%%%%%%%%%%%
% Ustawienie niestandardowej numeracji sekcji i podsekcji.
% Dodanie kropki po numerach (np. z "2.1 Rozdział" na "2.1. Rozdział").
% Można odkomentować, jeśli chcemy standardową formę numeracji - wtedy odkomentować wszystko poniżej do %%spis treści%%%%
\renewcommand{\thesection}{\arabic{chapter}.\arabic{section}.} % Numeracja sekcji w formacie: numer_rozdziału.numer_sekcji (np. 2.3.)
\renewcommand{\thesubsection}{\arabic{chapter}.\arabic{section}.\arabic{subsection}.} % Numeracja podsekcji w formacie: numer_rozdziału.numer_sekcji.numer_podsekcji (np. 2.3.1.)




%%%%%%%%%%%%%%%%%%%%%%%%Spis treści%%%%%%%%%%%%%%%%%%%%%%%%%%%%%
\setlength{\cftbeforechapskip}{0pt} %ustawia odstęp przed rozdziałami w spisie treści na 0pt; można zmienić wartość (np. 10pt), by zwiększyć odstęp, lub zakomentować tę linię, by przywrócić domyślne ustawienie

%%%%%%%%%%%%%%%%%%%%%%%%Tytuły rozdziałów%%%%%%%%%%%%%%%%%%%%%%%
\usepackage{titlesec}
\renewcommand{\cftchapaftersnum}{.}%ustawia, że po numerze rozdziału w spisie treści pojawi się kropka .
%\renewcommand{\chaptername}{}
%\renewcommand{\thechapter}{\arabic{chapter}.} % Dodaje kropkę po numerze rozdziału

\usepackage{ifthen}
\usepackage{ragged2e}  % do lewego wyrównania z lepszym łamaniem niż standardowy \raggedright

%instrukcja warunkowa pozwalająca na wybór formatowania nagłówka/tytułu rozdziału

\ifthenelse{\equal{\chapterstyletype}{classic}}{
  % Styl: Rozdział X (na górze), tytuł pod spodem
  \titleformat{\chapter}[display]
    {\normalfont\LARGE\bfseries\RaggedRight}
    {\chaptertitlename\ \thechapter} % np. "Rozdział 1"
    {0pt}
    {\LARGE}
}{%
  % Styl: X. Tytuł rozdziału (inline) z wiszącym wcięciem i wyrównaniem wielowierszowym
  \titleformat{\chapter}[block]
    {\normalfont\LARGE\bfseries\RaggedRight}
    {}
    {0pt}
    {%
      \makebox[1em][l]{\thechapter.}%
      \hangindent=1em
      \hangafter=1
      \ignorespaces
    }
}

% Styl rozdziału bez numeru (dla Bibliografii, Streszczenia itd.)
\titleformat{name=\chapter, numberless}[block]
  {\normalfont\LARGE\bfseries\RaggedRight}
  {}
  {0pt}
  {\hangindent=1em\hangafter=1}

%styl sekcji
\titleformat{\section}[hang]
  {\normalfont\Large\bfseries\RaggedRight} % format
  {\thesection}               % section number with dot
  {0.3em}                       % space between number and title
  {}                          % before title code
%styl podsekcji
  \titleformat{\subsection}[hang]
  {\normalfont\large\bfseries\RaggedRight} % format
  {\thesubsection}               % section number with dot
  {0.3em}                       % space between number and title
  {}                          % before title code
%stylpodpodsekcji
  \titleformat{\subsubsection}[hang]
  {\normalfont\normalsize\bfseries} % format
  {}               % brak numeru, czyli pusty label
  {0pt}            % odstęp między label a tytułem (brak, bo label pusty)
  {}                % before title code

% Ustawienie odstępów wokół tytułu rozdziału:
% \titlespacing*{<komenda>}{<lewy margines>}{<przed tytułem>}{<po tytule>}
% Tutaj:
% - lewy margines       = 0pt (brak wcięcia)
% - odstęp przed tytułem = 0pt (brak pustego miejsca na górze)
% - odstęp po tytule     = \baselineskip (dokładnie jeden wiersz odstępu przed tekstem)
\titlespacing*{\chapter}{0pt}{0pt}{\baselineskip}


%definicje, twierdzenia,...
\newtheorem{twierdzenie}{Twierdzenie}[chapter]
\newtheorem{lemat}[twierdzenie]{Lemat}
\newtheorem{wniosek}[twierdzenie]{Wniosek}
\newtheorem{obserwacja}[twierdzenie]{Obserwacja}
\newtheorem{uwaga}[twierdzenie]{Uwaga}
\theoremstyle{definition} %prostowanie czcionki
\newtheorem{przyklad}[twierdzenie]{Przykład}
\newtheorem{definicja}[twierdzenie]{Definicja}


%używane pakiety
%\usepackage{setspace}
%\usepackage{cite}
\usepackage{tikz}
\usetikzlibrary{positioning,arrows.meta,fit,backgrounds,shapes.geometric}
%\usepackage{subcaption}


\usepackage{algorithm}
\usepackage{algpseudocode}
\usepackage{listings}






\definecolor{niebieski}{RGB}{0,130,150} %kolor niebieski
\definecolor{mygray}{gray}{0.6} 		% kolor szary
\definecolor{zielony}{rgb}{0.98, 1.0, 0.98} % kolor zielony

\lstset{ 
	backgroundcolor=\color{zielony},   % choose the background color; you must add \usepackage{color} or \usepackage{xcolor}; should come as last argument
	basicstyle=\sffamily, %\footnotesize\ttfamily\fontsize{11}{0.8em}\selectfont,        % the size of the fonts that are used for the code
	breakatwhitespace=false,         % sets if automatic breaks should only happen at whitespace
	breaklines=true,                 % sets automatic line breaking
	captionpos=b,                    % sets the caption-position to bottom
	columns=fullflexible,
	commentstyle=\color{mygreen},    % comment style
	escapeinside={\%*}{*)},          % if you want to add LaTeX within your code
	extendedchars=true,              % lets you use non-ASCII characters; for 8-bits encodings only, does not work with UTF-8
	firstnumber=1,                   % start line enumeration with line 1
	frame=single,	                   % adds a frame around the code
	keepspaces=true,                 % keeps spaces in text, useful for keeping indentation of code (possibly needs columns=flexible)
	keywordstyle=\color{blue},       % keyword style
	numbers=left,                    % where to put the line-numbers; possible values are (none, left, right)
	numbersep=5pt,                   % how far the line-numbers are from the code
	numberstyle=\tiny\color{mygray}, % the style that is used for the line-numbers
	rulecolor=\color{black},         % if not set, the frame-color may be changed on line-breaks within not-black text (e.g. comments (green here))
	xleftmargin=0pt,%\parindent,
	showspaces=false,                % show spaces everywhere adding particular underscores; it overrides 'showstringspaces'
	showstringspaces=false,          % underline spaces within strings only
	showtabs=false,                  % show tabs within strings adding particular underscores
	language=python,
	morekeywords={*, end, function, ...}
	stepnumber=1,                    % the step between two line-numbers. If it's 1, each line will be numbered
	stringstyle=\color{mymauve},     % string literal style
	tabsize=2	                   % sets default tabsize to 2 spaces
}



%parametry dla kodów źródłowych
\definecolor{mygreen}{rgb}{0,0.6,0}
\definecolor{mygray}{rgb}{0.5,0.5,0.5}
\definecolor{mymauve}{rgb}{0.58,0,0.82}




\begin{document}
    \thispagestyle{empty}
\linespread{1.0}
\newgeometry{top=2.5cm, bottom=2.5cm, left=3cm, right=3cm}

	\begin{titlepage}
		\def\dystans{1cm}
		\begin{center}
			\Large 
\begin{figure}[h]
    \centering
    \includegraphics[height=3.2cm]{prz_pl.png}%
    \hfill
    \includegraphics[height=3.2cm]{wmifs_pl.png}
\end{figure} 
			{\large\scshape Politechnika Rzeszowska}
            
            {\large im. Ignacego Łukasiewicza}
        
            {\large\scshape\MakeUppercase{Wydział Matematyki i Fizyki Stosowanej}}
		\end{center}
		\vspace{1.3cm}
\begin{center}
  \begin{minipage}[c][5cm][c]{10cm}
    \centering
    {\Large Mikołaj Data} \par\vspace{0.4cm}
    {\Large\textbf{\textsc{Modelowanie dużych sieci LoRa Mesh}}} \par\vspace{0.4cm}
    Promotor: dr Krzysztof Pupka
  \end{minipage}
\end{center}
        
% \begin{center}
%     \begin{minipage}[c][5cm][c]{10cm} % 10 cm szerokości, 5 cm wysokości, wyrównane centralnie
%         \begin{center}
%             {\Large Imię i nazwisko autora}
%         \end{center}    
%         \begin{center}
%             {\Large\textbf{\textsc{Segmentacja zbioru sonogramów stawów biodrowych niemowląt ukierunkowana pod automatyzację oceny morfologii stawu biodrowego}}}
%         \end{center}
%         \begin{center}
%             Promotor: tytuł/stopień naukowy, imię i nazwisko
%         \end{center}
%     \end{minipage}
% \end{center}

\vspace{2cm}

        
		\begin{center}
			{\large  \textsc{Praca  inżynierska} \\
				\vspace{0.5cm}
				\large kierunek studiów:  Inżynieria i analiza danych}
		\end{center}
		\vspace{1.1cm}
		
		

		
		
		\vfill 
		\begin{center}
			{\large Rzeszów 2026}
		\end{center}
	\end{titlepage}
    \restoregeometry %lub strona_tytulowa-2
    \linespread{1.3}
    \newpage 
%----------------------------------------------------------------------------------
%                       PUSTA STRONA PO STRONIE TYTUŁOWEJ
%----------------------------------------------------------------------------------


    \thispagestyle{empty} %pusta strona bez numeru po stronie tytułowej
    ~\vfill
	\newpage
    
%---------------------------------------------------------------------------------
%                      SPIS TREŚCI
%---------------------------------------------------------------------------------
	\renewcommand{\contentsname}{\LARGE Spis treści}
    \tableofcontents
    \setcounter{tocdepth}{1}
    %\setcounter{page}{1} %numeracja stron od 1
	%\thispagestyle{headings}
    \cleardoublepage %wymusza rozpoczęcie nowego rozdziału na stronie nieparzystej (prawej); dodaje pustą stronę, jeśli potrzeba
    %\clearpage  %kończy bieżącą stronę i wypisuje wszystkie zaległe elementy (np. rysunki, tabele), ale nie wymusza rozpoczęcia nowego rozdziału na stronie nieparzystej
    %\newpage %Wymusza przejście do nowej strony, ale nie kończy otwartych elementów

%---------------------------------------------------------------------------------
%                      WSTĘP
%---------------------------------------------------------------------------------
	
	\chapter*{Wstęp} 
    \addcontentsline{toc}{chapter}{Wstęp}
	%\thispagestyle{empty}%strona nie jest numerowana
Technologie radiowe rozwijają się w dynamicznym tempie, a wraz z nimi pojawiają się nowe rozwiązania wykorzystywane do budowy bezprzewodowych sieci komunikacyjnych. Jednym z przykładów jest technologia LoRa, wykorzystująca modulację o tej samej nazwie, która została wprowadzona na rynek w 2009 roku. Jej charakterystyczne cechy, takie jak duży zasięg transmisji przy niewielkim zużyciu energii, sprawiły, że znalazła zastosowanie między innymi w systemach Internetu Rzeczy oraz rozproszonych sieciach radiowych.

W ostatnich latach rozwijane są również rozwiązania wykorzystujące LoRa do tworzenia rozproszonych sieci typu mesh, w których poszczególne urządzenia mogą pełnić jednocześnie funkcję nadawczą, odbiorczą oraz pośredniczyć w przekazywaniu danych pomiędzy innymi węzłami. Rozwiązania tego typu zaczęły zyskiwać na znaczeniu około 2020 roku, wraz z pojawieniem się pierwszych projektów umożliwiających budowę praktycznych sieci LoRa mesh.

Dynamiczny rozwój tego rodzaju sieci powoduje potrzebę opracowania narzędzi pozwalających na analizę ich funkcjonowania w różnych warunkach. Szczególnie istotna jest możliwość badania wpływu liczby i rozmieszczenia węzłów, parametrów transmisji oraz topologii sieci na jej wydajność i niezawodność. W tym celu przydatne są środowiska symulacyjne umożliwiające odwzorowanie zachowania sieci przed jej rzeczywistym wdrożeniem. Dostępność dedykowanych narzędzi pozwalających na symulację specyficznych mechanizmów współczesnych sieci LoRa mesh jest jednak ograniczona, co stanowi podstawę do podjęcia niniejszego tematu.
	
    \cleardoublepage %

%%%%%%%%%%%%%%%%%%%%% --- Zasadnicza treść rozprawy zaczyna się tutaj --- %%%%%%%%%%%%%%%%%%%%%

        \chapter{Podstawy transmisji radiowej LoRa}

\section{Modulacja LoRa}
Jest to autorska technika opatentowana przez firmę Semtech, stosowana w komunikacji radiowej. Wykorzystują ją urządzenia o niskim poborze mocy, często zasilane bateryjnie. LoRa korzysta z pasm ISM (przemysłowych, naukowych i medycznych), które nie wymagają licencji – w Europie 433,05–434,79 MHz oraz 863–870 MHz, a w Ameryce Północnej 902–928 MHz. Meshtastic w regionie EU\_868 korzysta z pasma 869,4--869,65 MHz. Technologia ta wywodzi się z metody modulacji CSS (ang. \textit{Chirp Spread Spectrum}), w której szybkość transmisji danych jest obniżana w celu poprawy czułości. Zgodnie z twierdzeniem Shannona-Hartleya, zwiększenie szerokości pasma poprawia parametry połączenia w odniesieniu do stosunku sygnału do szumu. W systemach wykorzystujących rozpraszanie widma (ang. \textit{Spread Spectrum}) przesyłane dane są mnożone przez sekwencję chipów (zazwyczaj sygnał o częstotliwości znacznie wyższej niż częstotliwość sygnału informacyjnego). Po stronie odbiorczej sygnał jest mnożony przez tę samą sekwencję chipów, co pozwala na odzyskanie oryginalnych danych. W przeciwieństwie do standardowej komunikacji z rozpraszaniem widma, LoRa wykorzystuje sygnał typu „chirp”, czyli sygnał, którego częstotliwość zmienia się w czasie.

\begin{itemize}
    \item \textbf{Częstotliwość środkowa} – częstotliwość, wokół której transmitowany jest sygnał.
    \item \textbf{Szerokość pasma} – w technice radiowej (RF) jest to zakres częstotliwości (mierzony w hercach) zajmowany przez sygnał.
    \item \textbf{Współczynnik rozpraszania (ang. \textit{Spreading Factor})} – w typowych zastosowaniach przyjmuje wartości od 7 do 12, a nowsze układy SX126x obsługują również wartości 5 i 6; parametr ten wpływa na szybkość transmisji danych, czas zajętości kanału (airtime), czas pracy na baterii oraz czułość odbiornika.
    \item \textbf{Współczynnik kodowania (ang. \textit{Code Rate})} – określa udział bitów niosących użyteczną informację do powtórzonych bitów używanych podczas korekcji błędów.
\end{itemize}

\section{Wpływ doboru parametrów na szybkość transmisji}

Jednym z istotnych zagadnień podczas projektowania sieci LoRa jest dobór parametrów transmisji, ponieważ mają one bezpośredni wpływ na szybkość przesyłania danych, czas zajęcia kanału oraz zasięg łączności. Do najważniejszych parametrów należą szerokość pasma (\(BW\)), współczynnik rozpraszania (\(SF\)) oraz współczynnik kodowania (\(CR\)).

Szybkość symbolowa w modulacji LoRa określana jest zależnością:

	\begin{equation}\label{R_S}
        R_S = \frac{BW}{2^{SF}}
    \end{equation}

gdzie \(R_S\) oznacza szybkość transmisji symboli w symbolach na sekundę, natomiast \(BW\) oznacza szerokość pasma sygnału. Zależność ta pokazuje, że zwiększenie szerokości pasma powoduje proporcjonalny wzrost szybkości symbolowej. Z kolei zwiększenie współczynnika rozpraszania powoduje wykładniczy spadek szybkości symbolowej.

Czas trwania pojedynczego symbolu można wyznaczyć jako odwrotność szybkości symbolowej:

	\begin{equation}\label{T_S}
        T_S = \frac{1}{R_S} = \frac{2^{SF}}{BW}
    \end{equation}

Wzrost wartości \(SF\) powoduje więc wydłużenie czasu transmisji pojedynczego symbolu. Jest to jedna z podstawowych zależności wpływających na charakterystykę transmisji LoRa. Wyższe wartości \(SF\) zwiększają odporność transmisji i umożliwiają uzyskanie większego zasięgu, jednak odbywa się to kosztem szybkości transmisji oraz zwiększenia czasu zajęcia kanału.

Przybliżoną szybkość transmisji danych można określić za pomocą zależności:

	\begin{equation}\label{R_b}
        R_b = SF \cdot \frac{BW}{2^{SF}} \cdot CR
    \end{equation}

gdzie \(R_b\) oznacza szybkość transmisji bitów, a \(CR\) współczynnik kodowania stosowany przez mechanizm korekcji błędów. Wartość \(CR\) może przyjmować wartości od \(4/5\) do \(4/8\). Zwiększenie wartości \(CR\) oznacza zastosowanie większej ilości nadmiarowych danych przeznaczonych do korekcji błędów, co poprawia odporność transmisji, lecz jednocześnie zmniejsza efektywną szybkość przesyłania danych.

W modulacji LoRa każdy symbol jest reprezentowany przez \(2^{SF}\) chipów. Szybkość transmisji chipów określana jest zależnością:

	\begin{equation}\label{R_C}
        R_C = R_S \cdot 2^{SF}
    \end{equation}

Po podstawieniu równania opisującego szybkość symbolową otrzymujemy:

	\begin{equation}\label{wzór}
        R_C = \frac{BW}{2^{SF}}\cdot 2^{SF}=BW
    \end{equation}

Oznacza to, że szybkość chipowa jest bezpośrednio równa szerokości pasma. Zwiększenie \(SF\) nie zmienia zatem szybkości generowania chipów, lecz powoduje zwiększenie liczby chipów przypadających na pojedynczy symbol, a w konsekwencji wydłużenie czasu jego transmisji.

Powyższe zależności pokazują występowanie kompromisu pomiędzy szybkością transmisji a parametrami odpowiedzialnymi za odporność i zasięg transmisji. Zwiększenie \(SF\) pozwala na uzyskanie większego zasięgu kosztem przepływności i czasu zajęcia kanału, natomiast zwiększenie \(BW\) pozwala zwiększyć szybkość transmisji kosztem zmniejszenia odporności na wąskopasmowe zakłócenia. Dobór parametrów transmisji powinien być zatem uzależniony od wymagań konkretnej sieci oraz warunków propagacyjnych.

\section{Budżet łącza radiowego oraz wpływ sprzętu}

Budżet łącza radiowego (ang. \textit{link budget}) jest zestawieniem wszystkich zysków i strat mocy występujących podczas transmisji sygnału pomiędzy nadajnikiem a odbiornikiem. Pozwala on oszacować moc sygnału docierającego do wejścia odbiornika oraz określić, czy przy założonych parametrach transmisja może zostać poprawnie odebrana.

Moc sygnału na wejściu odbiornika można w uproszczeniu określić za pomocą zależności:

	\begin{equation}\label{linkbudget}
        P_{RX} = P_{TX} + G_{TX} - L_{TX} - L_{FS} - L_M + G_{RX} - L_{RX}
    \end{equation}

gdzie:

\begin{itemize}
\item \(P_{RX}\) -- moc sygnału na wejściu odbiornika [dBm],
\item \(P_{TX}\) -- moc wyjściowa nadajnika [dBm],
\item \(G_{TX}\) -- zysk anteny nadawczej [dBi],
\item \(L_{TX}\) -- straty w torze nadawczym [dB],
\item \(L_{FS}\) -- tłumienie propagacyjne w wolnej przestrzeni (ang. \textit{Free Space Path Loss}) [dB],
\item \(L_M\) -- dodatkowe straty propagacyjne [dB],
\item \(G_{RX}\) -- zysk anteny odbiorczej [dBi],
\item \(L_{RX}\) -- straty w torze odbiorczym [dB].
\end{itemize}

Tłumienie w wolnej przestrzeni można wyznaczyć na podstawie częstotliwości oraz odległości pomiędzy nadajnikiem i odbiornikiem:

	\begin{equation}\label{L_FS}
        L_{FS} = 20\log_{10}(d) + 20\log_{10}(f) + 32{,}44
    \end{equation}

gdzie \(d\) oznacza odległość pomiędzy urządzeniami wyrażoną w kilometrach, a \(f\) częstotliwość w megahercach.

W rzeczywistym środowisku propagacja sygnału jest bardziej złożona niż model wolnej przestrzeni. Na poziom odbieranego sygnału wpływają między innymi przeszkody terenowe i zabudowa, odbicia, wielodrogowość, warunki atmosferyczne oraz niedopasowanie polaryzacji anten. Z tego powodu w praktycznych obliczeniach należy uwzględnić dodatkowy margines lub straty propagacyjne oznaczone jako \(L_M\).

Istotnym elementem budżetu łącza jest również zastosowany sprzęt. Zwiększenie mocy nadajnika bezpośrednio zwiększa moc sygnału docierającego do odbiornika, jednak jego wpływ na zasięg nie jest liniowy. Przykładowo zwiększenie mocy nadajnika z 1 W do 2 W oznacza wzrost mocy o około 3 dB, a nie dwukrotne zwiększenie zasięgu. Z punktu widzenia budżetu łącza taki wzrost odpowiada jedynie dodaniu 3 dB do bilansu energetycznego.

Duże znaczenie mają również anteny. Ich zysk wpływa bezpośrednio na budżet łącza, przy czym istotne jest nie tylko zastosowanie anteny o określonym zysku, ale również jej prawidłowe umiejscowienie, orientacja oraz zgodność polaryzacji anten nadawczej i odbiorczej. Straty występujące w przewodach koncentrycznych, złączach oraz innych elementach toru radiowego zmniejszają natomiast moc sygnału, która może zostać wypromieniowana lub odebrana.

W przypadku sieci LoRa szczególne znaczenie ma także czułość odbiornika. O poprawności transmisji nie decyduje wyłącznie moc odebranego sygnału \(P_{RX}\), lecz również jego relacja względem minimalnego poziomu sygnału wymaganego przez odbiornik dla zastosowanych parametrów modulacji. Różnica pomiędzy mocą odebranego sygnału a czułością odbiornika stanowi margines łącza:

	\begin{equation}\label{M}
        M = P_{RX} - S_{RX}
    \end{equation}

gdzie \(M\) oznacza margines łącza [dB], a \(S_{RX}\) czułość odbiornika [dBm]. Dodatni margines oznacza, że poziom odebranego sygnału znajduje się powyżej deklarowanej czułości odbiornika. Im większy jest margines, tym większa jest odporność połączenia na dodatkowe straty wynikające ze zmiennych warunków propagacyjnych.

Właściwy dobór elementów sprzętowych może zatem znacząco wpłynąć na rzeczywisty zasięg sieci LoRa. Zastosowanie anten o odpowiednim zysku, ograniczenie strat w torze radiowym oraz wykorzystanie urządzeń o wysokiej czułości odbiornika może przynieść większą poprawę parametrów łącza niż samo zwiększenie mocy nadajnika. Jest to szczególnie istotne w sieciach mesh, w których jakość poszczególnych połączeń pomiędzy węzłami wpływa na możliwość przekazywania danych przez kolejne węzły.

\section{Modelowanie strat ścieżki}

Straty ścieżki (ang. \textit{path loss}) określają zmniejszenie mocy sygnału radiowego podczas jego propagacji pomiędzy nadajnikiem a odbiornikiem. Ich wartość zależy między innymi od częstotliwości, odległości pomiędzy urządzeniami, wysokości anten oraz charakterystyki środowiska, w którym odbywa się transmisja. Dokładne określenie strat propagacyjnych jest istotnym elementem modelowania sieci radiowych, ponieważ bezpośrednio wpływa na przewidywany poziom sygnału odbieranego przez poszczególne węzły.

Modele strat ścieżki można podzielić na trzy podstawowe grupy: modele statystyczne, deterministyczne oraz kombinatoryczne. Modele statystyczne opracowywane są na podstawie wyników pomiarów przeprowadzonych w określonym środowisku. Ich zaletą jest stosunkowo niewielka złożoność obliczeniowa, jednak ich dokładność może znacząco zmniejszać się w przypadku zastosowania w środowisku różniącym się od tego, dla którego zostały opracowane.

Modele deterministyczne wykorzystują fizyczne prawa opisujące propagację fal radiowych. Uwzględniają między innymi geometrię obszaru, przeszkody znajdujące się pomiędzy nadajnikiem i odbiornikiem oraz zjawiska takie jak odbicie, dyfrakcja czy rozpraszanie. Pozwalają dzięki temu dokładniej odwzorować rzeczywiste warunki propagacji, jednak zazwyczaj wymagają większej liczby danych wejściowych i większej mocy obliczeniowej.

Modele kombinatoryczne łączą elementy podejścia deterministycznego i statystycznego. Wykorzystują informacje dotyczące fizycznych warunków propagacji, takich jak profil terenu, a jednocześnie uwzględniają statystyczny charakter niektórych zjawisk występujących podczas transmisji. Pozwala to uzyskać kompromis pomiędzy dokładnością modelowania a złożonością obliczeniową, co jest szczególnie istotne w przypadku symulacji obejmujących dużą liczbę węzłów.

Do prostych modeli strat ścieżki należy model propagacji w wolnej przestrzeni. Zakłada on brak przeszkód pomiędzy nadajnikiem i odbiornikiem oraz brak wpływu odbić i innych zjawisk charakterystycznych dla rzeczywistego środowiska. Straty w wolnej przestrzeni można wyznaczyć za pomocą wzoru \eqref{L_FS} Model ten jest prosty i charakteryzuje się bardzo małą złożonością obliczeniową, jednak nie uwzględnia rzeczywistego ukształtowania terenu ani znajdujących się na drodze sygnału przeszkód.

Bardziej rozbudowane modele wykorzystują informacje o charakterystyce środowiska. Przykładem modeli statystycznych są modele Okumura-Hata oraz COST-231, które zostały opracowane na podstawie wyników pomiarów propagacji fal radiowych. Modele te mogą zapewniać dobre rezultaty w określonych typach środowisk, takich jak obszary miejskie, podmiejskie lub wiejskie, jednak ich zastosowanie jest ograniczone przez zakres częstotliwości, odległości oraz warunki, dla których zostały opracowane.

W przypadku symulacji sieci rozmieszczonych na dużym obszarze istotne staje się uwzględnienie rzeczywistego ukształtowania terenu. W takiej sytuacji zastosowanie znajduje model Longley-Rice, znany również jako \textit{Irregular Terrain Model} (ITM). Model został opracowany w celu prognozowania propagacji fal radiowych w środowisku o nieregularnym ukształtowaniu terenu. W odróżnieniu od prostych modeli empirycznych wykorzystuje informacje dotyczące profilu terenu pomiędzy nadajnikiem i odbiornikiem, a także parametry opisujące środowisko propagacyjne.

Model Longley-Rice uwzględnia kilka mechanizmów mających wpływ na propagację sygnału, w tym propagację w wolnej przestrzeni, dyfrakcję, odbicia od powierzchni Ziemi oraz wpływ krzywizny Ziemi. Model uwzględnia również właściwości elektryczne powierzchni oraz parametry atmosferyczne. W zależności od warunków propagacji możliwe jest uwzględnienie różnych mechanizmów dominujących na danym odcinku trasy pomiędzy nadajnikiem i odbiornikiem.

Istotną cechą modelu Longley-Rice jest możliwość wykorzystania rzeczywistego profilu terenu. Dla określonej pary nadajnik--odbiornik model może otrzymać informacje o wysokości terenu w kolejnych punktach znajdujących się pomiędzy obiema antenami. Pozwala to uwzględnić sytuacje, w których bezpośrednia droga pomiędzy węzłami jest częściowo lub całkowicie zasłonięta przez wzniesienia terenu. Jest to szczególnie istotne w przypadku sieci radiowych rozmieszczonych na dużym obszarze, gdzie różnice wysokości mogą znacząco wpływać na możliwość nawiązania łączności.

Do podstawowych danych wejściowych modelu Longley-Rice należą między innymi częstotliwość pracy, odległość pomiędzy nadajnikiem i odbiornikiem, wysokości anten oraz profil terenu. Dodatkowo model może uwzględniać parametry środowiska, takie jak przewodność i przenikalność elektryczna powierzchni, współczynnik załamania atmosferycznego oraz charakterystykę klimatyczną obszaru. Wynikiem obliczeń jest między innymi przewidywana wartość strat propagacyjnych, która może następnie zostać wykorzystana do wyznaczenia mocy sygnału docierającego do odbiornika.

W niniejszej pracy do modelowania strat propagacyjnych wykorzystany zostanie model Longley-Rice. Jego wybór wynika przede wszystkim z charakteru symulowanej sieci. Węzły sieci LoRa mogą być rozmieszczone na znacznym obszarze, a ich wzajemna widoczność radiowa może być silnie uzależniona od ukształtowania terenu. Zastosowanie modelu uwzględniającego profil terenu pozwala zatem uzyskać bardziej realistyczne wyniki niż wykorzystanie wyłącznie modelu propagacji w wolnej przestrzeni lub prostego modelu empirycznego.

Wybór modelu Longley-Rice pozwala również zachować rozsądny kompromis pomiędzy dokładnością odwzorowania warunków propagacyjnych a złożonością obliczeniową symulatora. Jest to szczególnie istotne w przypadku sieci mesh, w których liczba potencjalnych połączeń pomiędzy węzłami rośnie wraz z liczbą urządzeń. Dla każdej pary węzłów możliwe jest wyznaczenie strat propagacyjnych, a następnie, na podstawie mocy nadajnika, zysków anten oraz czułości odbiornika, określenie możliwości nawiązania połączenia radiowego.

Otrzymana wartość strat propagacyjnych może zostać następnie wykorzystana w obliczeniach budżetu łącza ze wzoru \eqref{linkbudget}. W pracy korzystam z już gotowej biblioteki itmlogic \cite{oughton2020}.

Do obliczeń używamy numerycznego modelu terenu Shuttle Radar Topography Mission (SRTM).

\cleardoublepage
\chapter{Modelowanie warstwy logicznej sieci LoRa Mesh}
    Sieci LoRa Mesh stanowią rodzaj bezprzewodowych sieci kratowych, w których urządzenia wykorzystujące technologię LoRa mogą nie tylko bezpośrednio komunikować się ze sobą, ale również przekazywać dane pomiędzy kolejnymi węzłami. Dzięki temu informacja może zostać dostarczona do urządzenia docelowego poprzez wiele pośrednich węzłów, co pozwala zwiększyć zasięg sieci bez konieczności stosowania centralnej infrastruktury.

    W tego typu sieciach szczególne znaczenie ma warstwa logiczna odpowiedzialna za sposób adresowania urządzeń, przekazywania pakietów oraz wyboru kolejnych węzłów na drodze do odbiorcy. Właściwe zaprojektowanie tych mechanizmów pozwala wykorzystać zalety technologii LoRa, takie jak duży zasięg i niewielkie zapotrzebowanie na energię, przy jednoczesnym tworzeniu rozproszonych i odpornych na awarie sieci komunikacyjnych.
    
    \section{Popularne rozwiązania LoRa Mesh}
    Rozwój praktycznych rozwiązań wykorzystujących sieci LoRa Mesh w ostatnich latach doprowadził do powstania kilku otwartych projektów, które umożliwiają budowę zdecentralizowanych sieci komunikacyjnych bez wykorzystania infrastruktury operatorów telekomunikacyjnych. Do najpopularniejszych współczesnych rozwiązań należą przede wszystkim Meshtastic oraz MeshCore.
    
    Meshtastic jest otwartym projektem rozwijanym z myślą o komunikacji poza zasięgiem tradycyjnych sieci telekomunikacyjnych. Wykorzystuje urządzenia wyposażone w moduły LoRa do tworzenia zdecentralizowanej sieci, w której węzły mogą przekazywać informacje pomiędzy sobą. System umożliwia między innymi przesyłanie krótkich wiadomości tekstowych, informacji o położeniu oraz danych telemetrycznych. Istotną cechą rozwiązania jest możliwość korzystania z sieci bez dostępu do Internetu, sieci komórkowej lub innych elementów infrastruktury zewnętrznej.
    
    Drugim rozwijanym rozwiązaniem jest MeshCore, które koncentruje się na budowie zdecentralizowanej infrastruktury komunikacyjnej wykorzystującej LoRa. W architekturze MeshCore występują między innymi urządzenia użytkowników oraz węzły pełniące funkcję przekaźników, umożliwiające zwiększenie zasięgu sieci poprzez wielokrotne przekazywanie danych. Projekt rozwija również rozwiązania pozwalające na pracę urządzeń jako samodzielnych komunikatorów oraz integrację z aplikacjami działającymi na komputerach i urządzeniach mobilnych.
    
    Oba rozwiązania wykorzystują tę samą podstawową technologię transmisji radiowej, jednak różnią się zastosowanymi mechanizmami warstwy logicznej, sposobem przekazywania pakietów oraz organizacją poszczególnych rodzajów węzłów. Niniejsza praca dotyczy sieci Meshtastic, która jest obecnie najbardziej rozpowszechnionym z tych rozwiązań i dla której dostępne są publiczne dane z działającej sieci, wykorzystane w rozdziale~\ref{chap:weryfikacja} do weryfikacji symulatora. W dalszej części pracy przedstawiono mechanizmy warstwy logicznej sieci Meshtastic.
    
    \section{Podstawowe zasady działania routingu w sieciach Meshtastic}

    Routing w sieci Meshtastic opiera się na zdecentralizowanym przekazywaniu pakietów pomiędzy znajdującymi się w swoim zasięgu węzłami. W przeciwieństwie do klasycznych sieci, w których stosuję się topologię gwiazdy (np. w LoraWAN), w Meshtastic poszczególne urządzenia samodzielnie podejmują decyzję o przekazaniu odebranego pakietu. Pozwala to na tworzenie sieci, których działanie nie jest uzależnione od pojedynczego punktu infrastruktury.
    
    Węzeł wysyłający wiadomość nie musi znajdować się bezpośrednio w zasięgu urządzenia docelowego. Pakiet może zostać przekazany przez kolejne węzły, dzięki czemu komunikacja jest możliwa również pomiędzy urządzeniami oddalonymi od siebie o więcej niż jeden skok (ang. \textit{hop}). Liczba możliwych przekazań pakietu jest ograniczana za pomocą mechanizmu TTL (ang. \textit{Time To Live}), który zapobiega nieograniczonemu krążeniu pakietów w sieci.
    
    Istotnym elementem routingu Meshtastic jest mechanizm zarządzania retransmisjami. Węzły mogą przekazywać odebrane pakiety dalej, jednak nie każdy odebrany pakiet musi zostać ponownie wysłany. Decyzja o retransmisji zależy między innymi od rodzaju pakietu, jego parametrów oraz informacji posiadanych przez dany węzeł. Wykorzystywane są również mechanizmy pozwalające ograniczyć wielokrotne przekazywanie tego samego pakietu przez różne węzły.
    
    Aby zoptymalizować przekazywanie pakietów i tworzenie infrastruktury do sieci Meshtastic wprowadzano mechanizm ról. Są one przypisywane poszczególnym węzłom i wpływają na sposób przekazywania pakietów.
\begin{longtable}{|p{3.6cm}|p{5.2cm}|p{5.2cm}|}
\caption{Porównanie ról w sieci Meshtastic \cite{meshtastic_device}}
    \label{tabela1}\\
    \hline
    \textbf{Rola} & \textbf{Sposób działania} & \textbf{Zastosowanie} \\
    \hline
    \endfirsthead
    \hline
    \textbf{Rola} & \textbf{Sposób działania} & \textbf{Zastosowanie} \\
    \hline
    \endhead
    CLIENT &
    Urządzenie z połączoną aplikacją lub samodzielne urządzenie do przesyłania wiadomości. Retransmituje pakiety, jeżeli wcześniej nie zrobił tego żaden inny węzeł. &
    Ogólne zastosowanie dla użytkowników komunikujących się za pośrednictwem sieci Meshtastic, z obsługą aplikacji klienckich. \\
    \hline
    CLIENT\_MUTE &
    Urządzenie nie przekazuje pakietów pochodzących od innych urządzeń. &
    Udział w sieci bez wspierania routingu pakietów, co pozwala ograniczyć obciążenie sieci. \\
    \hline
    CLIENT\_HIDDEN &
    Urządzenie nadaje pakiety tylko wtedy, gdy jest to konieczne, ograniczając swoją widoczność oraz zużycie energii. &
    Zastosowania wymagające ograniczenia widoczności urządzenia lub zużycia czasu antenowego i energii. \\
    \hline
    CLIENT\_BASE &
    Osobista stacja bazowa, która zawsze retransmituje pakiety pochodzące od lub przeznaczone dla obserwowanych węzłów. Pozostałe pakiety obsługuje jak urządzenie CLIENT. &
    Stacje bazowe umieszczone w korzystnych lokalizacjach, np. na poddaszu lub dachu, w celu zwiększenia zasięgu własnych urządzeń. \\
    \hline
    ROUTER &
    Węzeł infrastrukturalny zwiększający zasięg sieci poprzez jednokrotną retransmisję każdego pakietu. Jest widoczny na liście węzłów. &
    Umieszczany w strategicznych lokalizacjach w celu zwiększenia zasięgu sieci. Jest prezentowany w topologii. \\
    \hline
    ROUTER\_LATE &
    Węzeł infrastrukturalny, który retransmituje pakiet jednokrotnie, ale dopiero po pozostałych węzłach. Pozwala zapewnić dodatkowe pokrycie lokalnych skupisk urządzeń. &
    Stosowany do pokrywania obszarów pozbawionych zasięgu lub zwiększenia niezawodności komunikacji w lokalnych skupiskach węzłów. \\
    \hline
    REPEATER &
    Węzeł infrastrukturalny, który wyłącznie retransmituje pakiety i nie nadaje własnych informacji o węźle, pozycji ani telemetrii. Nie jest widoczny na liście węzłów. Rola oznaczona jako przestarzała w wersji 2.7.11 oprogramowania. &
    Zwiększanie zasięgu sieci przy minimalnym narzucie. Nadal spotykana w działającej sieci, w nowych instalacjach zastępowana rolą ROUTER. \\
    \hline
\end{longtable}

Oprócz ról przedstawionych w tabeli firmware udostępnia role przeznaczone dla urządzeń o określonym zastosowaniu: TRACKER i SENSOR, które priorytetowo nadają odpowiednio pozycję i telemetrię, TAK i TAK\_TRACKER do współpracy z aplikacją ATAK oraz LOST\_AND\_FOUND, okresowo rozgłaszającą pozycję urządzenia. Symulator pozwala przypisać węzłowi każdą z tych ról, ponieważ konfiguracja przekazywana jest bezpośrednio do oprogramowania \texttt{meshtasticd}. Wyjątkiem są przestarzałe role REPEATER oraz ROUTER\_CLIENT, które oprogramowanie w wersji 2.7.15 zapisuje jako CLIENT. Zostało to sprawdzone przez ustawienie każdej roli w węźle \texttt{meshtasticd} i odczytanie jej po ponownym uruchomieniu.

\section{Problemy związane z symulacją warstwy logicznej sieci LoRa Mesh}

Symulacja warstwy logicznej sieci LoRa Mesh wiąże się z koniecznością odwzorowania zarówno struktury sieci, jak i mechanizmów odpowiedzialnych za przekazywanie pakietów pomiędzy węzłami. Jednym z podstawowych problemów jest dynamiczny charakter sieci, w której dostępność poszczególnych połączeń zależy od zasięgu radiowego, położenia urządzeń oraz aktualnych warunków propagacyjnych.

Istotnym wyzwaniem jest również odwzorowanie mechanizmów routingu i retransmisji. W rzeczywistym systemie wiele węzłów może odebrać ten sam pakiet i podjąć decyzję o jego dalszym przekazaniu. Symulator musi zatem uwzględniać ograniczenia liczby retransmisji, czas życia pakietów oraz możliwość występowania duplikatów. Pominięcie tych mechanizmów może prowadzić do wyników znacząco odbiegających od zachowania rzeczywistej sieci.

Kolejnym problemem jest modelowanie współdzielenia kanału radiowego. Urządzenia LoRa korzystają ze wspólnego medium, dlatego jednoczesne transmisje mogą prowadzić do kolizji i utraty pakietów. Dokładne odwzorowanie tego zjawiska wymaga uwzględnienia parametrów transmisji, czasu zajęcia kanału oraz wzajemnego oddziaływania nadających węzłów.

Złożoność symulacji rośnie wraz z liczbą węzłów oraz liczbą możliwych ścieżek pomiędzy nimi. Z tego względu istotnym elementem projektowanego modelu jest zachowanie odpowiedniego kompromisu pomiędzy dokładnością odwzorowania rzeczywistej sieci a czasem wykonywania symulacji.

Ponadto sam protokół dynamicznie się zmienia co powoduje, że w realnych zastosowaniach węzły mogą mieć różne wersje standardu i sposobu działania co tworzy kolejne utrudnienia.

\section{Symulacja węzła z użyciem MeshtasticD}
MeshtasticD jest natywnym programem binarnym służącym do uruchamiania Meshtastic na urządzeniach wyposażonych w radiomodemy komunikujące się za pośrednictwem interfejsu SPI lub USB. Jest dostępny dla systemów macOS oraz Linux. Dzięki wykorzystaniu modułu simradio można uruchomić program bez potrzeby posiadania fizycznego radiomodułu, dzięki czemu symulację działania pojedynczego węzła można uruchomić w izolowanym środowisku np. z użyciem kontenera Docker.

\cleardoublepage
\chapter{Implementacja symulatora}

Symulator ma jedno główne założenie. Logika sieci ma być prawdziwa, a symulowane jest tylko radio. Każdy węzeł to działający program Meshtastic, ten sam, który pracuje na fizycznych urządzeniach. Zamiast anteny i modułu LoRa dostaje on jednak pośrednika, który decyduje, kto usłyszy jego transmisję i z jaką mocą. Dzięki temu nie trzeba przepisywać routingu Meshtastic od nowa. Wystarczy dobrze udawać eter.

\section{Architektura}

Całość działa w kontenerach Docker. Na rysunku~\ref{rys:architektura} widać, jak części łączą się ze sobą.

\begin{figure}[ht]
\centering
\small
\begin{verbatim}
 przeglądarka (React + Leaflet)
          |  REST
 backend (FastAPI)
   |- precompute ITM  ->  macierz łączy (link_matrix.npz)
   |- RadioGateway    <-> węzeł 1 .. węzeł N  (meshtasticd -s, TCP)
   |- generator ruchu ->  harmonogram YAML
   `- zapis śladu     ->  InfluxDB  ->  Grafana
\end{verbatim}
\caption{Schemat architektury symulatora}
\label{rys:architektura}
\end{figure}

Backend napisany jest w Pythonie, w oparciu o FastAPI. Przechowuje scenariusze, czyli opis środowiska radiowego i listę węzłów z ich położeniem, wysokością anteny, rolą oraz limitem skoków. Uruchamia też kontenery węzłów i pilnuje całego cyklu życia symulacji.

Każdy węzeł to osobny kontener z programem \texttt{meshtasticd} uruchomionym w trybie \texttt{-s}. W tym trybie Meshtastic nie szuka fizycznego radia, tylko korzysta z modułu SimRadio. Każda transmisja, którą węzeł chciałby nadać w eter, trafia wtedy do jego klienta TCP jako pakiet typu \texttt{SIMULATOR\_APP}. W drugą stronę działa to podobnie. Pakiet \texttt{SIMULATOR\_APP} wysłany do węzła jest przez niego traktowany jak odebrany z anteny. Ten mechanizm jest punktem zaczepienia dla całego symulatora.

Przed startem symulacji każdy węzeł jest konfigurowany przez jego API. Dostaje region EU\_868, preset modemu, limit skoków, rolę i stałą pozycję. Potem jest restartowany, a symulator sprawdza, czy konfiguracja przetrwała restart. Jeśli coś się nie zgadza, start zostaje przerwany. Lepiej zatrzymać się od razu niż liczyć wyniki na węźle, który pracuje z innym presetem niż zakłada model radiowy.

Interfejs użytkownika to aplikacja React z mapą Leaflet. Prowadzi przez kolejne kroki: ustawienie środowiska, rozmieszczenie węzłów na mapie, obliczenie łączy, uruchomienie symulacji i podgląd pakietów. Szczegółowy ślad trafia do bazy InfluxDB, a gotowe pulpity w Grafanie pokazują na przykład odsetek dostarczonych pakietów albo przebieg retransmisji konkretnej wiadomości.

\section{Obliczanie łączy przed symulacją}

Model Longley-Rice jest kosztowny, więc liczę go raz, przed uruchomieniem węzłów. Dla każdej pary węzłów symulator buduje profil terenu z danych SRTM \cite{farr2007}. Próbki biorę mniej więcej co 90 metrów, bo taka jest rozdzielczość danych SRTM3. Początkowo profil miał zawsze 50 odcinków. Dla trasy długości 40 km oznaczało to próbkę co 800 metrów, a przy takiej gęstości łatwo przeoczyć wzniesienie. Po zagęszczeniu profilu straty dla dłuższych tras zmieniły się nawet o 6 dB.

Na profilu liczona jest strata terenowa ITM z biblioteki itmlogic \cite{oughton2020}. Przyjmuję typowe parametry gleby (przenikalność 15, przewodność 0,005 S/m), refrakcję 314 jednostek N oraz medianę rozkładu, czyli niezawodność i ufność po 50\%. Dla obszarów, dla których są dostępne modele budynków LoD1 w formacie CityGML, symulator sprawdza dodatkowo, czy budynki wchodzą w pierwszą strefę Fresnela. Jeśli tak, dolicza stratę dyfrakcyjną dominującej przeszkody.

Wynikowa mediana straty ścieżki dla pary nadajnik i odbiornik ma postać:

\begin{equation}\label{eq:pl_sym}
L = L_{ITM} + L_{bud} + L_{ot} + X_{\ell}, \qquad X_{\ell} \sim \mathcal{N}(0, \sigma_{\ell}^2)
\end{equation}

gdzie \(L_{ITM}\) to strata terenowa, \(L_{bud}\) strata od budynków, \(L_{ot}\) stała strata otoczenia, a \(X_{\ell}\) trwałe cieniowanie łącza. Dwa ostatnie składniki wymagają krótkiego komentarza. Model terenu nie widzi drzew, gęstej zabudowy poza obszarem z danymi CityGML, anteny schowanej na parapecie ani kiepskiego kabla. Strata otoczenia zbiera te wszystkie rzeczy w jedną liczbę. Cieniowanie łącza robi coś innego. Losuje dla każdej pary węzłów jedną wartość, taką samą w obu kierunkach, i trzyma ją przez całą symulację. Dzięki temu dwa łącza tej samej długości nie muszą być tak samo dobre, co w prawdziwej sieci zdarza się bardzo często. Oba parametry zostały dobrane na danych z działającej sieci, o czym więcej w rozdziale~\ref{chap:weryfikacja}. Domyślnie są wyzerowane.

Z budżetu łącza \eqref{linkbudget} wynika medianowa moc odbierana:

\begin{equation}\label{eq:rssi}
P_{RX} = P_{TX} + G_{TX} + G_{RX} - L
\end{equation}

Domyślna moc nadajnika to 22 dBm. Tyle daje układ SX1262, który siedzi w większości popularnych urządzeń, a Meshtastic domyślnie nadaje z maksymalną mocą układu. Limit regionu wynosi co prawda 27 dBm, ale niewiele węzłów go osiąga.

Szum odbiornika i czułość wynikają z tego samego wzoru:

\begin{equation}\label{eq:szum}
N = -174 + 10\log_{10}(BW) + NF, \qquad S_{RX} = N + SNR_{min}(SF)
\end{equation}

gdzie \(BW\) jest szerokością pasma w hercach, \(NF = 6\) dB to współczynnik szumów odbiornika, a \(SNR_{min}\) minimalny stosunek sygnału do szumu, przy którym demodulator LoRa jeszcze sobie radzi. Dla pasma 250 kHz szum wynosi \(-114\) dBm. Dla presetu MEDIUM\_FAST (SF9, \(SNR_{min} = -12{,}5\) dB) daje to czułość \(-126{,}5\) dBm.

Wynikiem tego etapu jest macierz łączy zapisywana na dysku. Razem z nią zapisuję odcisk wszystkich danych wejściowych modelu. Gdy użytkownik zmieni coś, co wpływa na propagację, odcisk przestaje pasować i łącza trzeba policzyć od nowa. Zmiana nazwy węzła albo seeda fadingu macierzy nie unieważnia.

\section{Eter w czasie działania symulacji}

Za eter odpowiada moduł RadioGateway. Jest połączony z każdym węzłem i nasłuchuje pakietów \texttt{SIMULATOR\_APP}. Kiedy któryś węzeł nadaje, gateway wie dwie rzeczy: kto nadaje i kiedy zaczął. Nadawcą jest węzeł, którego połączenie zgłosiło pakiet. Pole \texttt{from} w pakiecie się do tego nie nadaje, bo przy retransmisji wskazuje autora wiadomości, a nie węzeł, który właśnie ją powtórzył.

Dla każdego sąsiada nadawcy gateway losuje fading \(X \sim \mathcal{N}(0, \sigma^2)\) i liczy moc odebraną jako \(P_{RX} + X\). Fading jest losowany osobno dla każdej transmisji. Opisuje więc to, że ten sam pakiet wysłany dwa razy po tej samej trasie raz dojdzie, a raz nie. SNR to różnica mocy i szumu, ale obcięta do 12 dB, bo tyle mniej więcej pokazuje prawdziwy odbiornik LoRa.

To obcięcie wygląda na kosmetykę, a okazało się bardzo ważne. Bez niego węzły stojące blisko siebie dostawały SNR rzędu 70 dB. Meshtastic wylicza z SNR opóźnienie retransmisji (im lepszy sygnał, tym dłużej węzeł czeka, żeby pierwszeństwo miały węzły dalsze). Przy 70 dB opóźnienie wyszło na około 45 minut. Pakiet stał w kolejce nadawczej, blokował wszystko za sobą i węzeł po prostu milkł. Z zewnątrz wyglądało to tak, jakby symulacja gubiła ruch. Dopiero log firmware z wpisem o opóźnieniu 2752920 ms pokazał, co się dzieje.

Przy okazji wyszedł drugi problem. Biblioteka meshtastic dla Pythona przed wysłaniem pakietu do węzła czeka, aż w jego kolejce nadawczej zwolni się miejsce. Czeka bez limitu czasu. Wszystkie odbiory były wstrzykiwane z jednego wątku tej biblioteki, więc jeden zapchany węzeł zatrzymywał eter dla całej sieci. Teraz gateway zapisuje odebrany pakiet bezpośrednio do strumienia TCP. To ma sens, bo z punktu widzenia węzła jest to odbiór z anteny, a nie wiadomość od aplikacji, więc kolejka nadawcza w ogóle go nie dotyczy.

\section{Kolizje}

Na początku symulator nie znał kolizji. Każdy pakiet powyżej czułości po prostu docierał. Dla kilku węzłów to nie przeszkadza. Ale sieć Meshtastic działa trochę jak sala, w której wszyscy mówią na jednym kanale. Przy kilkudziesięciu węzłach i zalewowym rozsyłaniu wiadomości to właśnie nakładanie się transmisji najbardziej ogranicza sieć. Bez kolizji wyniki skalowalności byłyby zdecydowanie zbyt optymistyczne.

Model kolizji opieram na pracy Bora i współautorów \cite{bor2016} i na regułach z ich symulatora LoRaSim. Pierwsza rzecz to czas trwania pakietu w eterze. Liczę go wzorem z noty aplikacyjnej Semtech \cite{semtech_an1200_13}:

\begin{equation}\label{eq:airtime}
\begin{split}
T_{pkt} &= \left(n_{pr} + 4{,}25\right) T_S + n_{dane}\, T_S, \\
n_{dane} &= 8 + \max\left(\left\lceil \frac{8PL - 4SF + 44}{4(SF - 2DE)} \right\rceil CR_d,\ 0\right)
\end{split}
\end{equation}

gdzie \(T_S\) to czas symbolu ze wzoru~\eqref{T_S}, \(n_{pr} = 16\) to długość preambuły używana przez Meshtastic, \(PL\) długość pakietu w bajtach (dane i 16 bajtów nagłówka), \(CR_d\) mianownik współczynnika kodowania (od 5 do 8), a \(DE\) przyjmuje wartość 1, gdy symbol trwa dłużej niż 16 ms. Dla presetu LONG\_FAST i pakietu 16 bajtów daje to 354 ms.

Odbiór oceniam dopiero wtedy, gdy pakiet się skończy. Prawdziwy odbiornik też oddaje pakiet do programu dopiero po odebraniu całości, więc gateway czeka czas \(T_{pkt}\) i jeszcze 50 ms zapasu. W tym momencie zna już wszystkie transmisje, które nałożyły się na dany pakiet, także te, które zaczęły się później. Z kodu SimRadio \cite{meshtastic_firmware} wynika, że w standardowej kompilacji \texttt{meshtasticd} przetwarza wstrzyknięty pakiet od razu, więc opóźnienie się nie podwaja.

Pakiet \(p\) nie dociera do odbiorcy \(r\), jeśli zachodzi któryś z warunków:
\begin{itemize}
\item jego moc u \(r\) jest poniżej czułości,
\item \(r\) sam nadawał w czasie trwania \(p\), bo radio LoRa nie odbiera podczas nadawania,
\item nakłada się z nim pakiet \(q\), którego moc u \(r\) spełnia warunek
\begin{equation}\label{eq:capture}
P_p - P_q < \Delta_{cap}, \qquad \Delta_{cap} = 6~\textrm{dB},
\end{equation}
chyba że \(q\) skończył się w pierwszych \(n_{pr} - 5 = 11\) symbolach preambuły pakietu \(p\).
\end{itemize}

Ostatni warunek to efekt przechwycenia. Jeśli jeden z dwóch pakietów jest wyraźnie silniejszy, odbiornik go wyłapie, a słabszy przepadnie. Gdy oba są podobnie silne, giną oba. Wyjątek dla preambuły też pochodzi z LoRaSim. Odbiornik do synchronizacji potrzebuje tylko kilku ostatnich symboli preambuły, więc zakłócenie, które skończy się wcześniej, nie robi mu krzywdy.

Zostaje jeszcze jedna rzecz. Prawdziwy węzeł Meshtastic nie nadaje na ślepo. Najpierw odczekuje losowy czas z okna rywalizacji, a potem sprawdza, czy ktoś już nadaje (mechanizm CAD, \textit{Channel Activity Detection}). Jeśli kanał jest zajęty, losuje opóźnienie od nowa \cite{meshtastic_firmware}. SimRadio losowe opóźnienie zachowuje, ale CAD pomija, bo dla niego kanał nigdy nie jest zajęty. Bez poprawki węzły, które się dobrze słyszą, wchodziłyby sobie w słowo znacznie częściej niż w rzeczywistości. Dlatego gateway emuluje CAD. Jeśli węzeł zaczyna nadawać, a w tym momencie słyszy cudzy sygnał powyżej czułości, jego transmisja jest przesuwana za koniec tamtego sygnału o losową liczbę slotów od 0 do 7. Długość slotu i rozmiar okna biorę z kodu firmware: slot to 2,5 czasu symbolu plus 7,6 ms.

Każda nieudana próba odbioru trafia do śladu razem z przyczyną: za słaby sygnał, kolizja albo nadawanie odbiorcy. W widoku pakietów i w Grafanie widać więc od razu, ile strat robią kolizje, a ile sam zasięg.

Model został sprawdzony na prawdziwych węzłach \texttt{meshtasticd} w klasycznym układzie ukrytego węzła. Węzły A i C stoją po dwóch stronach węzła B i nie słyszą się nawzajem. Ich sygnały u B różnią się o 5,7 dB, czyli mniej niż próg przechwycenia. A i C nadawały jednocześnie w ośmiu rundach. Z włączonymi kolizjami B stracił oba pakiety w sześciu rundach. W dwóch pozostałych losowe opóźnienie w firmware rozsunęło transmisje na tyle, że się nie nałożyły. Po wyłączeniu kolizji B odebrał wszystkie szesnaście pakietów.

Model ma swoje uproszczenia. Synchronizacja na preambule jest opisana jedną regułą jedenastu symboli. Wszystkie węzły pracują z tym samym współczynnikiem rozpraszania, więc nie modeluję częściowej ortogonalności różnych SF. Moc zakłócenia jest stała przez cały czas trwania pakietu. Emulowany CAD wykrywa sygnał przez cały czas transmisji, choć prawdziwy układ najpewniej łapie preambułę.

\section{Generowanie ruchu}

Ruch w symulacji opisuje plik YAML z harmonogramem. Każda pozycja mówi, kiedy i z którego węzła coś wysłać. Może to być wiadomość do wszystkich, wiadomość do jednego węzła albo traceroute. Generator przekazuje te polecenia węzłom tak, jak zrobiłaby to aplikacja na telefonie. Poza tym węzły nadają swój zwykły ruch, na przykład informacje o sobie i telemetrię. Symulator go nie wycina, bo w prawdziwej sieci ten ruch też zajmuje eter.

Fading i cieniowanie łączy mają własne seedy, więc da się powtórzyć przebieg z tym samym stanem eteru. Opóźnień losowanych wewnątrz \texttt{meshtasticd} symulator nie kontroluje, więc dwa przebiegi z tymi samymi seedami nie są identyczne co do milisekundy. Dlatego wyniki porównuję na kilku przebiegach, a nie na jednym.


\cleardoublepage
\chapter{Weryfikacja symulatora na danych z działającej sieci}
\label{chap:weryfikacja}

Symulator może działać bez zarzutu i nadal kłamać. Najprostszy sposób, żeby to wyłapać, to porównać go z prawdziwą siecią. W Polsce jest to możliwe dzięki serwisowi lorastats.pl \cite{lorastats}. Zbiera on pakiety polskiej sieci Meshtastic odebrane przez bramki MQTT, czyli węzły, które poza radiem mają też dostęp do internetu i przesyłają dalej wszystko, co usłyszą. Udostępnia je przez API z ostatnich 24 godzin.

\section{Dane}

Z API pobrane zostały cztery zbiory dla wszystkich regionów serwisu: listy węzłów, zdekodowane pakiety, surowe pakiety oraz raporty odbioru. Raport odbioru mówi, która bramka usłyszała który pakiet, z jaką mocą i SNR, ile skoków zostało pakietowi i który węzeł nadał go ostatni. Łącznie to około 1,07 miliona raportów od 360 bramek z 6 i 7 października 2026 roku.

Najwięcej kłopotu sprawiły pozycje węzłów. Meshtastic pozwala ograniczyć dokładność rozgłaszanej pozycji \cite{meshtastic_protobufs}. Węzeł ustawia liczbę bitów precyzji, a pozycja jest zaokrąglana do środka komórki siatki o boku \(2^{32-b} \cdot 10^{-7}\) stopnia. Przy domyślnych dla kanału publicznego 13 bitach komórka ma około 5,8 km. Z 1389 węzłów z pozycją tylko 258 podawało ją z pełną precyzją. Pozostałe są w pewnym sensie celowo rozmyte.

Jak bardzo to przeszkadza, pokazuje prosty eksperyment. Dla 25 par węzłów prawdziwe położenie losowane jest wewnątrz komórki siatki, a następnie porównywana jest strata ITM dla pozycji prawdziwej i dla tej, którą węzeł by rozgłosił. Wyniki są w tabeli~\ref{tab:precyzja}.

\begin{table}[ht]
\centering
\caption{Błąd straty ścieżki wynikający z ograniczonej precyzji pozycji}
\label{tab:precyzja}
\small
\begin{tabular}{|c|c|c|c|c|}
\hline
\textbf{Bity} & \textbf{Bok komórki} & \textbf{Błąd pozycji} & \textbf{Błąd straty} & \textbf{95. percentyl} \\
\hline
13 & 5,8 km & 1,9 km & 10,4 dB & 33,0 dB \\
15 & 1,5 km & 461 m & 5,4 dB & 19,4 dB \\
17 & 365 m & 119 m & 3,4 dB & 11,2 dB \\
19 & 91 m & 29 m & 1,2 dB & 4,1 dB \\
21 & 23 m & 7 m & 0,5 dB & 2,5 dB \\
\hline
\end{tabular}
\end{table}

Przy 13 bitach błąd jest większy niż cały fading. Takie węzły nie nadają się do sprawdzania modelu łączy. Dopiero od mniej więcej 19 bitów błąd pozycji przestaje mieć znaczenie. Przy okazji wyszedł błąd w symulatorze. Dwa węzły z zaokrągloną pozycją często lądują w tym samym punkcie, a dla odległości zero biblioteka itmlogic dzieliła przez zero i cały scenariusz nie dawał się policzyć. Teraz odległość jest ograniczona od dołu do jednego metra.

Do porównań wzięte zostały więc tylko węzły, które we wszystkich raportach podają pozycję z pełną precyzją, nie przemieszczają się i nie mają podejrzanie okrągłych współrzędnych. Zostało ich 149. Trzeba jednak pamiętać, że większość z nich ma pozycję wpisaną ręcznie. Pełna precyzja zapisu nie gwarantuje, że ktoś trafił w swój dach.

\section{Co można zobaczyć w raportach}

Bramka wysyła do serwera pierwszą kopię pakietu, jaką usłyszy. Oryginalna transmisja jest zawsze pierwsza w czasie. Jeśli więc bramka zgłasza pakiet bez zużytych skoków, a ostatnim nadawcą jest sam autor, to łącze bezpośrednie istnieje. Jeśli natomiast bramka słyszała dany węzeł wiele razy, ale zawsze dopiero po retransmisji, to bezpośredniego łącza najpewniej nie ma.

Nie każda bramka nadaje się na takiego obserwatora. Część z nich zgłasza każdy pakiet z wyzerowanym licznikiem skoków albo w ogóle nie raportuje odbiorów bezpośrednich. Wygląda to na mostki MQTT, a nie na radio. Takie bramki zostały pominięte, bo ich „brak łącza” nic nie znaczy.

Z węzłów o pełnej precyzji zbudowanych zostało osiem rozłącznych klastrów, po co najwyżej dziesięć węzłów, w różnych częściach kraju: okolice Zielonej Góry i Nowej Soli, Rzeszów, Lublin, Siedlce, Szczecin, Poznań i Warszawa. Preset każdego klastra wynika z hasha kanału w surowych pakietach. Większość pracuje na MEDIUM\_FAST, zachód kraju na SHORT\_SLOW, a Siedlce na LONG\_FAST. W klastrach jest 87 par nadajnik i bramka z wiarygodną obserwacją. W 31 z nich łącze bezpośrednie istnieje, w 56 nie.

\section{Łącza}

Pierwsza wersja modelu przewidywała łącze prawie wszędzie. Wykrywała 94\% prawdziwych łączy, ale poprawnie odrzucała tylko 36\% par, które w rzeczywistości się nie słyszą. Na parach odległych o 5 do 15 km model widział łącze w 18 przypadkach na 18. W rzeczywistości były cztery. Moc na działających łączach była zawyżona średnio o 13,8 dB.

Sama kolejność łączy była przy tym całkiem dobra. Pole pod krzywą ROC dla marginesu łącza wynosi 0,80. Model więc wie, które łącze jest lepsze, a które gorsze. Myli się co do tego, gdzie leży granica.

Pierwsze poprawki dotyczyły oczywistych rzeczy: mocy 22 dBm zamiast 27, szumu wyliczanego z pasma i gęstszego profilu terenu. Zawyżenie mocy spadło do 10,1 dB, a błąd SNR z 41 do 12 dB. To wciąż za mało.

Tu pojawia się ciekawa sprzeczność. Gdy dobiera się samą stałą stratę otoczenia, moc na działających łączach zgadza się najlepiej przy 5 do 10 dB. Odsetek istniejących łączy zgadza się jednak dopiero przy 20 do 30 dB. Rozwiązaniem jest trwałe cieniowanie łącza ze wzoru~\eqref{eq:pl_sym}. Oglądamy tylko te łącza, które działają, a działają głównie te, którym cieniowanie akurat sprzyja. Ich moc jest więc wyższa od mediany modelu i ta selekcja tłumaczy całą rozbieżność. Tabela~\ref{tab:cieniowanie} pokazuje to na liczbach.

\begin{table}[ht]
\centering
\caption{Zgodność modelu z danymi dla różnych wartości straty otoczenia i cieniowania}
\label{tab:cieniowanie}
\small
\begin{tabular}{|c|c|c|c|c|}
\hline
\(L_{ot}\) [dB] & \(\sigma_{\ell}\) [dB] & \textbf{Błąd mocy (dane)} & \textbf{Błąd mocy (model)} & \textbf{Łącza (model)} \\
\hline
10 & 10 & \(+0{,}1\) dB & \(-4{,}9\) dB & 0,58 \\
20 & 15 & \(-9{,}9\) dB & \(-10{,}2\) dB & 0,46 \\
25 & 20 & \(-14{,}9\) dB & \(-15{,}0\) dB & 0,41 \\
30 & 20 & \(-19{,}9\) dB & \(-17{,}5\) dB & 0,36 \\
\hline
\end{tabular}
\end{table}

W danych łącza istnieją w 36\% par. Strata otoczenia od 25 do 30 dB i cieniowanie z odchyleniem około 20 dB oddają jednocześnie i ten odsetek, i moc działających łączy. Dwadzieścia decybeli to dużo więcej niż typowe cieniowanie terenowe, które w literaturze wynosi zwykle od 6 do 12 dB. Ten parametr zbiera jednak wszystko, czego model nie wie o konkretnych instalacjach: prawdziwą wysokość anteny (znaną tylko dla 40 z 72 węzłów w klastrach), dokładność ręcznie wpisanej pozycji, kable i to, czy antena stoi na dachu, czy w pokoju. Część tej zmienności to cecha węzła, a nie łącza. Do opisania jej osobno dla każdego węzła danych jest jednak za mało.

\section{Routing}

Ostatni test jest najbliższy temu, do czego symulator powstał. Pięć klastrów zostało uruchomionych na prawdziwych węzłach \texttt{meshtasticd}, z rolami i limitami skoków odczytanymi z sieci. Każdy węzeł wysłał trzy wiadomości rozgłoszeniowe, a potem traceroute do głównej bramki klastra. Dla każdej wiadomości liczy się pierwszy udany odbiór na bramce: ile skoków zużył pakiet i który węzeł go ostatnio powtórzył. To samo da się odczytać z raportów lorastats.

Porównanie dotyczy 72 par nadawca i bramka oraz trzech wariantów modelu łączy. Wariant fizyczny to sam teren i budynki. Dwa warianty skalibrowane mają stratę otoczenia 27 dB, cieniowanie 20 dB i różnią się tylko seedem cieniowania. Wyniki są w tabeli~\ref{tab:routing}. Te przebiegi odbyły się jeszcze bez modelu kolizji.

\begin{table}[ht]
\centering
\caption{Routing w symulacji w porównaniu z siecią rzeczywistą}
\label{tab:routing}
\begin{tabular}{|p{6.6cm}|c|c|c|}
\hline
 & \textbf{Fizyczny} & \textbf{Seed 1} & \textbf{Seed 2} \\
\hline
Pary osiągnięte w symulacji & 51 & 51 & 44 \\
Łącza bezpośrednie (symulacja / sieć) & 44 / 12 & 25 / 11 & 24 / 11 \\
Zgodność: bezpośrednio czy przez pośrednika & 37\% & 65\% & 61\% \\
Ta sama liczba skoków & 27\% & 33\% & 36\% \\
Średnia różnica liczby skoków & \(-1{,}47\) & \(-0{,}92\) & \(-0{,}93\) \\
Ten sam węzeł pośredniczący z klastra & brak przypadków & 6 z 6 & 3 z 6 \\
\hline
\end{tabular}
\end{table}

W wariancie fizycznym prawie każdy węzeł słyszy bramkę bezpośrednio. Trasy wychodzą więc o półtora skoku krótsze niż w sieci, a retransmisji praktycznie nie ma czego porównywać. Po kalibracji zgodność co do tego, czy pakiet idzie bezpośrednio, czy przez pośrednika, rośnie z 37\% do ponad 60\%. Gdy w prawdziwej sieci pakiet przekazuje węzeł z klastra, symulacja wskazuje ten sam węzeł w dziewięciu przypadkach na dwanaście.

Część rozbieżności zostanie zawsze. W prawdziwej sieci wiele tras przechodzi przez węzły spoza klastra, które nie podają dokładnej pozycji i nie mogły zostać umieszczone w symulacji. Stąd w symulacji nadal brakuje średnio około jednego skoku, a mniej więcej co trzecia para jest w ogóle nieosiągalna. Widać też, że wynik zależy od losowania cieniowania. Wnioski o routingu trzeba więc wyciągać z kilku przebiegów.

Trzeba też uczciwie zaznaczyć jedną rzecz. Parametry kalibracji zostały dobrane na tych samych parach węzłów, na których potem sprawdzany jest routing. Kalibracja dotyczyła wyłącznie łączy, a routing robi niezmieniony firmware, więc test nie jest pozbawiony sensu. Nie jest jednak w pełni niezależny. Lepszą weryfikację dałby drugi, osobny zbiór danych, na przykład z innego okresu albo z pomiarów prowadzonych celowo.


    \cleardoublepage
    


    %-----------------------------------------------------------------------------------
    %                            BIBLIOGRAFIA
    %-----------------------------------------------------------------------------------
    \begin{thebibliography}{3}
    \addcontentsline{toc}{chapter}{Bibliografia}
    \bibitem{oughton2020} E. Oughton, T. Russell, J. Johnson, C. Yardim, J. Kusuma, „itmlogic: The Irregular Terrain Model by Longley and Rice”, \emph{Journal of Open Source Software}, vol. 5, nr 2266, 2020, doi: 10.21105/joss.02266.
    \bibitem{meshtastic_device} Meshtastic, \emph{Device Configuration}. Dostępne online: \url{https://meshtastic.org/docs/configuration/radio/device/} [dostęp: 03.09.2026].
    \bibitem{bor2016} M. Bor, U. Roedig, T. Voigt, J. M. Alonso, „Do LoRa Low-Power Wide-Area Networks Scale?”, w: \emph{Proceedings of the 19th ACM International Conference on Modeling, Analysis and Simulation of Wireless and Mobile Systems (MSWiM '16)}, ACM, 2016, doi: 10.1145/2988287.2989163.
    \bibitem{semtech_an1200_13} Semtech Corporation, \emph{LoRa Modem Designer's Guide}, nota aplikacyjna AN1200.13, wersja 1, lipiec 2013.
    \bibitem{farr2007} T. G. Farr i in., „The Shuttle Radar Topography Mission”, \emph{Reviews of Geophysics}, vol. 45, nr 2, RG2004, 2007, doi: 10.1029/2005RG000183.
    \bibitem{lorastats} lorastats.pl, \emph{Statystyki polskiej sieci Meshtastic}, API danych z ostatnich 24 godzin. Dostępne online: \url{https://lorastats.pl/API} [dostęp: 07.10.2026].
    \bibitem{meshtastic_protobufs} Meshtastic, \emph{Protobufs}, plik \texttt{meshtastic/mesh.proto}, pole \texttt{Position.precision\_bits}. Dostępne online: \url{https://github.com/meshtastic/protobufs} [dostęp: 07.10.2026].
    \bibitem{gugik_lod1} Główny Urząd Geodezji i Kartografii, \emph{Modele 3D budynków LoD1}, dane otwarte, licencja CC BY 4.0. Dostępne online: \url{https://opendata.geoportal.gov.pl/InneDane/Budynki3D/LOD1/2024/} [dostęp: 07.10.2026].
    \bibitem{itu_p526} International Telecommunication Union, \emph{Recommendation ITU-R P.526-15: Propagation by diffraction}, Genewa, 2019.
    \bibitem{meshtastic_firmware} Meshtastic, \emph{Meshtastic Firmware}, wersja 2.7.15, pliki \texttt{src/platform/portduino/SimRadio.cpp}, \texttt{src/mesh/RadioInterface.cpp}, \texttt{src/mesh/SX126xInterface.cpp}. Dostępne online: \url{https://github.com/meshtastic/firmware} [dostęp: 07.10.2026].
	\end{thebibliography}
    \cleardoublepage

    %-----------------------------------------------------------------------------------
    %                            SPISY RYSUNKÓW/TABEL
    %-----------------------------------------------------------------------------------
    %\renewcommand{\listfigurename}{\LARGE Spis rysunków}
    %\addcontentsline{toc}{chapter}{Spis rysunków}
    %\listoffigures
    %\cleardoublepage

    %\renewcommand{\listtablename}{\LARGE Spis tabel}
    %\addcontentsline{toc}{chapter}{Spis tabel}
    %\listoftables
    %\cleardoublepage

    %-----------------------------------------------------------------------------------
    %                            DODATKI/ANEKSY
    %-----------------------------------------------------------------------------------
	

    

    %-----------------------------------------------------------------------------------
    %                            STRESZCZENIA
    %-----------------------------------------------------------------------------------
	

    \cleardoublepage
	\thispagestyle{empty}
\setlength{\parindent}{0in}
\setlength{\parskip}{0em}
\linespread{1.0}\selectfont

\textbf{POLITECHNIKA RZESZOWSKA} \hfill Rzeszów 2026 \\
\textbf{im. Ignacego Łukasiewicza} \\
\textbf{WYDZIAŁ MATEMATYKI I FIZYKI STOSOWANEJ} \\

\begin{center}
	\textbf{STRESZCZENIE  PRACY  DYPLOMOWEJ }
\end{center}

\textbf{Tytuł:} Modelowanie dużych sieci LoRa Mesh \\
\textbf{Autor:} Mikołaj Data \\
\textbf{Promotor:} dr Krzysztof Pupka \\
\textbf{Słowa kluczowe:} LoRa, Meshtastic, sieci kratowe, symulacja, model Longley-Rice, routing
\\

Celem pracy było opracowanie symulatora sieci LoRa Mesh opartych na projekcie Meshtastic. Warstwa logiczna sieci, w tym routing i retransmisje, realizowana jest przez rzeczywiste oprogramowanie \texttt{meshtasticd} uruchamiane w kontenerze Docker dla każdego węzła, a symulowana jest wyłącznie warstwa radiowa. Straty propagacyjne wyznaczane są modelem Longley-Rice na podstawie danych SRTM, uzupełnionym o stratę od budynków, cieniowanie i fading. Odwzorowano również kolizje, półdupleks i wykrywanie zajętości kanału. Model łączy skalibrowano na danych z polskiej sieci Meshtastic (lorastats.pl) i oceniono na niezależnym zbiorze testowym: AUC wyniosło 0,79, a tryb przekazania pakietów odwzorowano w 75\% par węzłów. Trasy w symulacji są średnio o 1,5 skoku krótsze niż w sieci rzeczywistej. Działanie symulatora sprawdzono dla sieci liczących do 80 węzłów. Nadaje się on przede wszystkim do porównywania wariantów konfiguracji sieci.


% Streszczenia PL i EN nie mieszczą się na jednej stronie, więc EN zaczyna się od nowej
\clearpage
\thispagestyle{empty}



\textbf{RZESZOW UNIVERSITY OF TECHNOLOGY} \hfill Rzeszów 2026 \\
\textbf{FACULTY OF MATHEMATICS AND APPLIED PHYSICS } \\

\begin{center}
	\textbf{DIPLOMA THESIS ABSTRACT}
\end{center}

\textbf{Title:} Modelling large LoRa Mesh networks\\
\textbf{Author:} Mikołaj Data \\
\textbf{Supervisor:} Krzysztof Pupka, PhD \\ %% Angielski zapis stopni/tytułów naukowych promotora 

\textbf{Key words:} LoRa, Meshtastic, mesh networks, simulation, Longley-Rice model, routing
\\

The aim of the thesis was to develop a simulator of LoRa Mesh networks based on the Meshtastic project. The logical layer of the network, including routing and retransmissions, is handled by the actual \texttt{meshtasticd} firmware running in a Docker container for each node, while only the radio layer is simulated. Path loss is computed with the Longley-Rice model from SRTM data, extended with building loss, shadowing and fading. Collisions, half-duplex operation and channel activity detection are also modelled. The link model was calibrated on data from the Polish Meshtastic network (lorastats.pl) and evaluated on an independent test set: it achieved an AUC of 0.79, and the forwarding mode was reproduced for 75\% of node pairs. Routes in the simulation are on average 1.5 hops shorter than in the real network. The simulator was tested on networks of up to 80 nodes. It is best suited to comparing network configurations.

\vfill

% Poniżej przykłady angielskich odpowiedników polskich stopni i tytułów naukowych promotora:
% - doktor | dr → Dr Firstname Lastname lub Firstname Lastname, PhD
% - dr inż. → Dr Eng. Firstname Lastname lub Firstname Lastname, PhD Eng.
% - doktor habilitowany | dr hab. → Firstname Lastname, PhD, DSc
% - dr hab. inż. → Firstname Lastname, PhD, DSc Eng.
% - dr hab. prof. PRz → Firstname Lastname, PhD, DSc, Assoc. Prof.
% - prof. dr hab. → Prof. Firstname Lastname
% - prof. dr hab. inż. → Prof. Eng. Firstname Lastname
% - doktor, profesor PRz → Doctor, Firstname Lastname, Associate Professor lub Firstname Lastname, PhD, Assoc. Prof.

    \thispagestyle{empty}

    %-----------------------------------------------------------------------------------
    %                            TABELA GENAI
    %-----------------------------------------------------------------------------------
    \cleardoublepage
	
\thispagestyle{empty}
\begin{center}
    Wykaz obszarów i narzędzi GenAI \\
    wykorzystanych przez autora w trakcie wykonywania pracy dyplomowej
\end{center}

\begin{table}[ht]
    \centering
    \begin{tabularx}{\linewidth}{|c|X|X|}\hline
       \textbf{Lp.}  & \multicolumn{1}{c|}{\textbf{Obszar wykorzystywania $^{a)}$}} & \multicolumn{1}{c|}{\textbf{Narzędzie}} \\ \hline
        1. & a) redakcja i korekta tekstu & Claude (Anthropic) \\ \hline
        2. & h) pisanie i testowanie kodu & Claude (Anthropic) \\ \hline
    \end{tabularx}
    \label{tab:my_label}
\end{table}

\textbf{Uwaga!} W przypadku, gdy autor nie korzysta z obszarów i narzędzi GenAI w tabeli wpisuje „nie korzystano”

\

\textit{$^{a)}$ Przykładowe obszary wykorzystania:
    \begin{enumerate}[label=\alph*)]
        \item redakcja i korekta tekstu (m.in. poprawa stylistyczna i gramatyczna, synonimy i~parafrazowanie, sprawdzenie spójności językowej),
        \item analiza stanu wiedzy, przegląd literatury,
        \item generowanie treści lub przykładów,
        \item tworzenie pytań do ankiet i kwestionariuszy,
        \item tworzenie schematów, diagramów i map myśli,
        \item analiza danych i wizualizacja wyników (m.in. obliczenia statystyczne, tworzenie wykresów i tabel),
        \item tworzenie podsumowań i wniosków,
        \item pisanie i testowanie kodu,
        \item tłumaczenia językowe.
    \end{enumerate}} %obowiązkowe dla studentów wybierających tematy od semestru letniego 2024/2025, pozostali mogą usunąć
	
	\end{document}

	
